\section{OpenAI、医薬化学ベンチマークMedChemBenchを公開}
OpenAIは医薬化学の推論能力を測るMedChemBenchを発表した。117件の論文・特許を基にした99問で、分子構造の読み取り、実験結果の解釈、次に試験する化合物の選択などを評価する。公式資料ではGPT-6.2 Astraのxhigh設定が84.5\%を記録したと報告。薬剤候補探索の支援につながる可能性を示す一方、ベンチマーク成績と実際の創薬成果は区別して捉える必要がある。
\dailyxsource{OpenAI公式投稿}{https://x.com/OpenAI/status/2123856316187975756}

\section{OpenAI、モデル蒸留への対策研究を共有}
OpenAIは、安全対策を引き継がずにモデルの能力を移転する蒸留行為について、検出・阻止方法を含む研究を共有した。業界全体のセーフガード強化に資するとの説明だ。ただし、Grokが確認した投稿には論文やブログへの直接リンクがなく、具体的な検出方式や有効性の詳細は未確認である。モデル能力の再利用と安全管理の関係が焦点となる。
\dailyxsource{OpenAI公式投稿}{https://x.com/OpenAI/status/2123862344422088794}

\section{GPT-Rosalindによる薬剤発見の手掛かり}
OpenAIは、GPT-Rosalindが神経変性疾患に関する薬剤発見の新たな手掛かりを示したと投稿した。Perry症候群に関係するDCTN1遺伝子の作用機序について提案し、エージェント群が32分で文献とデータを結び付けたという。投稿時刻は観測窓の終端である07:00 JST。引用元や関連研究は窓外の可能性があり、提案の実験的検証は今後の確認事項となる。
\dailyxsource{OpenAI公式投稿}{https://x.com/OpenAI/status/2123899406550528422}

\section{Gemini 4 Neon、スキーシミュレータを生成}
Google DeepMindはGemini 4 Neonによるスキーシミュレータのデモを公開した。モンブランの地形を基に、カスタムシェーダーを用いたリアルタイム描画を実現したと説明。事前に用意した3Dモデルを使わずに構築したという。モデル自体の発表は観測窓外の可能性があるが、デモ投稿は窓内であり、コード生成が視覚的なアプリケーション制作に直結する例として注目された。
\dailyxsource{Google DeepMind公式投稿}{https://x.com/GoogleDeepMind/status/2123826572385833119}

\section{Anthropic、共同創業者の歩みを動画で紹介}
Anthropicは共同創業者たちが会社の歩みを語る動画を投稿した。IPOを巡る議論が続く中、企業の成り立ちや方向性に関心が集まっている。GrokによればIPO関連の公式声明は観測窓外の11月19日に出ており、今回の動画自体は窓内に投稿された。動画の詳しい内容やIPO条件については本入力だけでは確認できないため、確定的な説明は避ける。
\dailyxsource{Anthropic公式投稿}{https://x.com/AnthropicAI/status/2123822989154369958}

\section{Liquid AIのd1決定モデルに再び言及}
Liquid AIのd1モデルファミリーに関する投稿が観測された。エッジ機器での単パス推論や、画像・音声を扱う能力が特徴として紹介されている。関連するブログは以前のリリースである可能性があり、今回の投稿を新製品発表と同一視することはできない。低遅延の判断処理を端末側で行う構成が、エージェント実装の選択肢として関心を集めている。
\dailyxsource{Liquid AI関連投稿}{https://x.com/liquidai/status/2123861236714799147}

\section{Qwen4-Maxの提供状況を巡る未確認情報}
AlibabaのQwen4-Maxがプラットフォーム上で利用可能になったとの二次投稿が見られた。100万トークンの文脈長、深い推論、マルチモーダル対応などが主張されている。ただし、Grokは公式一次資料を確認できておらず、投稿時刻も07:07 JSTと観測窓の終了後である。提供開始や仕様を確定事実として扱わず、窓外の参考情報として記録する。
\dailyxsource{二次投稿・観測窓外}{https://x.com/francis_cgl/status/2123901072381288675}

\section{Genie 4世界モデルへの関心}
Google DeepMindのGenie 4を、リアルタイムで操作できるシミュレーションの研究プレビューとして紹介する投稿があった。信頼されたテスターへの提供が始まったとされるが、Grokが確認したのは二次情報であり、公式の一次投稿は未確認である。世界モデルを仮想環境の生成やエージェントの行動評価に用いる方向性への期待が読み取れる。
\dailyxsource{二次投稿}{https://x.com/industrial_XR/status/2123808733272698952}

\section{自律エージェントによるモデル組織運営の事例}
Hugging Face上のオープンモデル組織を自律エージェントが8週間運営したというケーススタディが紹介された。投稿では600時間超のGPU利用、38モデルの作成などが挙げられた。研究論文のURLはGrok入力で確認されておらず、数字も投稿者の主張として扱う。モデル開発から公開までの一連の作業をエージェントが担う際の実現性と監督方法が論点となる。
\dailyxsource{研究紹介投稿}{https://x.com/paperswithcodee/status/2123808907327885531}

\section{GLM-5.5の来週公開説}
GLM-5.5が翌週に公開されるとの噂が流れた。1兆パラメータ超、100万トークンの文脈長、オープンウェイト提供などが取り沙汰されているが、いずれも二次投稿に基づき、Zhipu AIによる公式確認はない。競合モデルとの性能比較も現段階では予測にすぎない。今後の公式発表を待つべき情報として整理する。
\dailyxsource{未確認の二次投稿}{https://x.com/theaidesk_/status/2123851117272408402}
